\documentclass[a4paper,11pt]{article}
% Shared preamble — % Shared preamble — % Shared preamble — \input{../shared/preamble} in every document
% Each document must declare \documentclass[a4paper,11pt]{article} before \input-ing this file.
\usepackage[T1]{fontenc}
\usepackage[utf8]{inputenc}
\usepackage[polish,english]{babel}
\usepackage[top=2cm,bottom=2cm,left=2cm,right=2cm]{geometry}
\usepackage{booktabs}
\usepackage{array}
\usepackage{tabularx}
\usepackage{longtable}
\usepackage{multirow}
\usepackage{hhline}
\usepackage{enumitem}
\usepackage{xcolor}
\usepackage{fancyhdr}
\usepackage{graphicx}
\usepackage{lastpage}

\definecolor{pzzblue}{RGB}{0,70,127}

\pagestyle{fancy}
\fancyhf{}
\renewcommand{\headrulewidth}{0.4pt}
\fancyfoot[C]{\thepage/\pageref{LastPage}}

% Helper: underlined fill-in field of given width
\newcommand{\field}[1]{\underline{\hspace{#1}}}
\newcommand{\fieldline}{\noindent\rule{\linewidth}{0.4pt}}
 in every document
% Each document must declare \documentclass[a4paper,11pt]{article} before \input-ing this file.
\usepackage[T1]{fontenc}
\usepackage[utf8]{inputenc}
\usepackage[polish,english]{babel}
\usepackage[top=2cm,bottom=2cm,left=2cm,right=2cm]{geometry}
\usepackage{booktabs}
\usepackage{array}
\usepackage{tabularx}
\usepackage{longtable}
\usepackage{multirow}
\usepackage{hhline}
\usepackage{enumitem}
\usepackage{xcolor}
\usepackage{fancyhdr}
\usepackage{graphicx}
\usepackage{lastpage}

\definecolor{pzzblue}{RGB}{0,70,127}

\pagestyle{fancy}
\fancyhf{}
\renewcommand{\headrulewidth}{0.4pt}
\fancyfoot[C]{\thepage/\pageref{LastPage}}

% Helper: underlined fill-in field of given width
\newcommand{\field}[1]{\underline{\hspace{#1}}}
\newcommand{\fieldline}{\noindent\rule{\linewidth}{0.4pt}}
 in every document
% Each document must declare \documentclass[a4paper,11pt]{article} before \input-ing this file.
\usepackage[T1]{fontenc}
\usepackage[utf8]{inputenc}
\usepackage[polish,english]{babel}
\usepackage[top=2cm,bottom=2cm,left=2cm,right=2cm]{geometry}
\usepackage{booktabs}
\usepackage{array}
\usepackage{tabularx}
\usepackage{longtable}
\usepackage{multirow}
\usepackage{hhline}
\usepackage{enumitem}
\usepackage{xcolor}
\usepackage{fancyhdr}
\usepackage{graphicx}
\usepackage{lastpage}

\definecolor{pzzblue}{RGB}{0,70,127}

\pagestyle{fancy}
\fancyhf{}
\renewcommand{\headrulewidth}{0.4pt}
\fancyfoot[C]{\thepage/\pageref{LastPage}}

% Helper: underlined fill-in field of given width
\newcommand{\field}[1]{\underline{\hspace{#1}}}
\newcommand{\fieldline}{\noindent\rule{\linewidth}{0.4pt}}


\fancyhead[L]{\textbf{SAFETY BRIEFING / ODPRAWA BEZPIECZEŃSTWA} — Bałtyk}
\fancyhead[R]{\today}

\begin{document}
\selectlanguage{polish}

\begin{center}
    {\Large\bfseries\color{pzzblue} ODPRAWA BEZPIECZEŃSTWA}\\
    {\large SAFETY BRIEFING}\\[4pt]
    {\normalsize Morze Bałtyckie \textbullet{} jacht: \field{5cm} \textbullet{} data: \field{3cm}}
\end{center}

\vspace{0.3cm}
\noindent\rule{\linewidth}{1pt}

% --- Jacht i wyposażenie ---
\section*{1. Jacht i wyposażenie / Vessel \& Equipment}

\begin{itemize}[noitemsep]
    \item Gaśnice: \field{5cm} (lokalizacja / location)
    \item Tratwa ratunkowa: \field{5cm} — system uruchomienia: \field{3cm}
    \item Koło ratunkowe: \field{4cm} — z linką świetlną: Tak / Nie
    \item Kamizelki ratunkowe: szafa \field{3cm} — każdy jest zobowiązany nosić w złych warunkach
    \item Uprzęże / liny asekuracyjne: \field{5cm}
    \item EPIRB: \field{5cm} — aktywacja automatyczna / ręczna
    \item Flary / rakiety: \field{5cm}
    \item Apteczka: \field{4cm}
\end{itemize}

% --- Procedury alarmowe ---
\section*{2. Procedury alarmowe / Emergency Procedures}

\subsection*{Człowiek za burtą (MOB)}
\begin{enumerate}[noitemsep]
    \item Krzyknąć \textbf{„CZŁOWIEK ZA BURTĄ!"} i wcisnąć MOB na GPS/ploterze.
    \item Rzucić koło ratunkowe / boję MOB.
    \item Wyznaczyć obserwatora — nie spuszczać wzroku z osoby w wodzie.
    \item Zawołać pozostałych — kapitan przejmuje stery.
    \item Podejście do osoby w wodzie — wybór manewru (Quick Stop / Reach\&Return).
    \item Radio VHF Ch~16: MAYDAY jeśli sytuacja krytyczna.
\end{enumerate}

\subsection*{Pożar}
\begin{enumerate}[noitemsep]
    \item Krzyknąć \textbf{„POŻAR!"} — odciąć gaz / paliwo.
    \item Użyć gaśnicy. Nie otwierać luku przed ugaszeniem.
    \item Przy braku opanowania: wysłać sygnał MAYDAY, gotowość do opuszczenia jachtu.
    \item Punkt zbiórki: \field{5cm}
\end{enumerate}

\subsection*{Przeciek / zalanie}
\begin{enumerate}[noitemsep]
    \item Zlokalizować źródło, zamknąć zawory denne.
    \item Uruchomić pompę zęzową: elektryczna \field{3cm} / ręczna \field{3cm}.
    \item Tamowanie szmatami, klinami.
    \item Przy zagrożeniu: MAYDAY, przygotować tratę ratunkową.
\end{enumerate}

% --- Komunikacja ---
\section*{3. Komunikacja radiowa / Radio Communication}

\begin{itemize}[noitemsep]
    \item Radio VHF: \field{3cm} — kanał wywołania i niebezpieczeństwa: \textbf{Ch~16}
    \item MRCC Gdynia: VHF Ch~16, tel.\ +48 58 620 62 00
    \item Prognoza pogody IMGW: VHF Ch~\field{1.5cm} (godz.\ \field{2cm})
    \item Kanał roboczy załogi: Ch~\field{1.5cm}
    \item MAYDAY: \textit{„MAYDAY MAYDAY MAYDAY, tu jacht [nazwa], pozycja [pozycja], [opis sytuacji], na pokładzie [liczba osób], wzywam pomocy, odbiór."}
\end{itemize}

% --- Bałtyk specyfika ---
\section*{4. Specyfika Bałtyku / Baltic-specific Hazards}

\begin{itemize}[noitemsep]
    \item Szybko zmieniająca się pogoda — obserwuj barometr i IMGW.
    \item Wody płytkie (Mierzeja Wiślana, Hel) — monitoruj głębokościomierz.
    \item Duży ruch statków handlowych — COLREGS, AIS.
    \item Sieci rybackie — wachta nocna, zwolnij w rejonach połowów.
    \item Zasolenie niskie — tratwa może zachowywać się inaczej niż na oceanie.
\end{itemize}

% --- Zasady ogólne ---
\section*{5. Zasady ogólne / General Rules}

\begin{itemize}[noitemsep]
    \item Nie wychodzić na pokład bez powiadomienia wachtowego (nocą: obowiązkowo uprząż).
    \item Zakaz spożywania alkoholu podczas wachty.
    \item Każdy zna lokalizację wyłącznika głównego silnika: \field{4cm}
    \item Każdy zna lokalizację zaworów dennych: \field{5cm}
    \item Obowiązek meldowania wszelkich nieprawidłowości kapitanowi.
\end{itemize}

\vspace{0.6cm}
\noindent\rule{\linewidth}{0.4pt}
\vspace{0.2cm}

\noindent\textbf{Potwierdzenie uczestnictwa w odprawie / Briefing acknowledgement:}

\vspace{0.3cm}
\begin{tabularx}{\linewidth}{|X|X|c|}
\hline
\textbf{Imię i nazwisko} & \textbf{Data} & \textbf{Podpis} \\
\hline
 & & \\[5pt]\hline
 & & \\[5pt]\hline
 & & \\[5pt]\hline
 & & \\[5pt]\hline
 & & \\[5pt]\hline
 & & \\[5pt]\hline
 & & \\[5pt]\hline
 & & \\[5pt]\hline
\end{tabularx}

\vspace{0.5cm}
\noindent Kapitan: \field{6cm} \hfill Podpis: \field{5cm}

\end{document}
